% =====================================================================
% Trajectory judge prompts.
%
% AppWorld trajectory judge
%   Source: data_synthesis/llm_judge/evaluation_prompts.py
%   Functions:
%     - get_appworld_solution_verifier_system_prompt
%     - get_appworld_solution_verifier_user_prompt
%
% OfficeBench trajectory judge
%   Source: data_synthesis/officebench/simulation/judge_prompts.py
%   Functions:
%     - get_officebench_solution_verifier_system_prompt
%     - get_officebench_solution_verifier_user_prompt
%
% Both prompts use the `prettyprompt' lstlisting style. Per-listing
% `captionpos=b' overrides the style's default top-positioned caption
% so captions sit below each listing. Each listing is wrapped in
% figure* to span both columns. The Unicode arrow character (->) in the
% AppWorld system prompt is shown as ASCII `->' to render reliably in
% \ttfamily without an extra literate mapping.
% =====================================================================


% =====================================================================
% AppWorld trajectory judge
% =====================================================================
\paragraph{AppWorld trajectory judge prompt.}
\Cref{lst:judge_appworld_system,lst:judge_appworld_user} show the prompts used by the LLM judge that filters AppWorld trajectories. The system prompt (\Cref{lst:judge_appworld_system}) defines nine requirements (documentation use, context understanding, strategy, information flow, parameter grounding, iteration, completion via \texttt{apis.supervisor.complete\_task}, deduplication, relevance) that a passing trajectory must satisfy. The user prompt (\Cref{lst:judge_appworld_user}) wraps the task and the agent's interaction trace in \texttt{<task>} / \texttt{<trajectory>} tags and asks for a YES/NO verdict on a line beginning with \texttt{`ANSWER:'}.


% --- AppWorld judge system prompt ---
\begin{figure*}[t]
\begin{lstlisting}[style=prettyprompt, captionpos=b, caption={AppWorld trajectory judge -- system prompt. Defines nine requirements that a passing trajectory must satisfy.}, label={lst:judge_appworld_system}]
You are an AI model with strong reasoning and code understanding capabilities. You are provided with a user task within <task></task> XML tags and an agent's solution trajectory for this task within <trajectory></trajectory> XML tags. The provided trajectory consists of multiple steps, each containing agent's reasoning, code generated by the agent that includes one or more API calls and the corresponding code execution output.

Your job is to analyze whether the agent's solution trajectory correctly and completely solves the given task. DO NOT use your own factual world knowledge when verifying. The code execution output may contain synthetic data - focus on whether the agent's logic and approach reflects how the given task should be solved.

The solution MUST satisfy these requirements:

1. Always follows the documentation:
   - An API is never called without querying its documentation first.

2. Understands full context of the task:
   - If the task refers to prior activity, the agent should retrieve relevant data to fully understand the user's intent.

3. Follows correct strategy:
   - The agent breaks down the task into appropriate steps and executes them in logical order.

4. Proper information flow:
   - Retrieves necessary information through API calls
   - Uses outputs from earlier code in subsequent code (e.g., using returned access_token, IDs, etc.)
   - Properly chains API calls (login -> authenticated requests)

5. Grounded API parameters:
   - API input parameters either come from the task description, or the API descriptions, or the code from previous steps, or the previous code execution outputs, and must match the API documentation

6. Handles iteration correctly:
   - When the task requires searching or aggregating all possible entities of a particular type, the agent's code should iterate through all of them
   - Correctly implements loops or pagination with appropriate page limit value (based on the documentation) to access all data when required
   - Properly tracks state values across iterations (max value, matching items, etc.)
   - When an API returns a list, the agent should go through it to correctly identify the target item(s) instead of assuming that the target items are at the beginning of the list

7. Reaches completion:
   - If the task asks for some information (e.g., 'How many songs?', 'What is the email?'), the value of the answer provided to `apis.supervisor.complete_task(answer=...)` should answer the user's question correctly and completely based on the retrieved data.
   - If the task does not require an answer, `apis.supervisor.complete_task()` must be called without passing any argument value.

8. Avoids duplication:
   - The agent does not make redundant API calls that retrieve the same data multiple times unnecessarily.

9. Avoids irrelevant actions:
   - The agent does not perform actions that are completely irrelevant to solving the task.

Minor inefficiencies are acceptable (e.g., retrieving slightly more data than needed).
\end{lstlisting}
\end{figure*}


% --- AppWorld judge user prompt ---
\begin{figure*}[t]
\begin{lstlisting}[style=prettyprompt, captionpos=b, caption={AppWorld trajectory judge -- user prompt. The agent's full interaction is inserted at \texttt{\{interaction\_trace\}}.}, label={lst:judge_appworld_user}]
Here is the task:
<task>
My name is {first_name} {last_name}. My personal email is {email} and phone number is {phone_number}.
{instruction}
</task>

Here is the agent's trajectory:
<trajectory>
{interaction_trace}
</trajectory>

Analyze whether the agent's solution trajectory correctly and completely solves the given task. First, provide your step-by-step thought process that clearly explains in detail why you think the agent's solution is either correct or incorrect. After providing your explanation, give your final answer in a separate line starting with `ANSWER:`. Your final answer must be either YES or NO. DO NOT provide anything else after the final YES/NO answer.
\end{lstlisting}
\end{figure*}


% =====================================================================
% OfficeBench trajectory judge
% =====================================================================
\paragraph{OfficeBench trajectory judge prompt.}
\Cref{lst:judge_officebench_system,lst:judge_officebench_user} show the prompts used by the LLM judge that filters OfficeBench trajectories. The OfficeBench system prompt (\Cref{lst:judge_officebench_system}) is rewritten from the AppWorld version to match OfficeBench's action format: actions are JSON \texttt{<answer>}-blobs rather than Python API calls, the agent must \texttt{switch\_app} before issuing app-specific actions, completion is signalled by \texttt{system.finish\_task} (not \texttt{apis.supervisor.complete\_task}), and many tasks are graded by side effects on files under \texttt{tasks/<task\_id>/testbed/data/}. The user prompt (\Cref{lst:judge_officebench_user}) follows the same \texttt{<task>}/\texttt{<trajectory>} structure but uses first-name only (OfficeBench supervisors expose nothing else).


% --- OfficeBench judge system prompt ---
\begin{figure*}[t]
\begin{lstlisting}[style=prettyprompt, captionpos=b, basicstyle=\ttfamily\scriptsize, caption={OfficeBench trajectory judge -- system prompt. Defines eight requirements tailored to OfficeBench's JSON-action interface and file-side-effect grading.}, label={lst:judge_officebench_system}]
You are an AI model with strong reasoning and code understanding capabilities. You are provided with a user task within <task></task> XML tags and an agent's solution trajectory for that task within <trajectory></trajectory> XML tags. The trajectory consists of multiple steps; each step contains the agent's JSON-formatted action wrapped as `<answer> {...} </answer>` and the observation returned by the OfficeBench environment.

OfficeBench is a multi-app office assistant: the agent can switch between apps (`calculator`, `calendar`, `email`, `excel`, `pdf`, `word`, `shell`, `llm`) and the meta `system` app. To use an app's action the agent must first issue `{"app":"system","action":"switch_app","target_app":"<app>"}`. The agent terminates the task with `{"app":"system","action":"finish_task","answer":<value>}`, where `<value>` is the literal answer for question-style tasks or the string `"None"` for tasks that don't ask for an answer. Files for the task live under `tasks/<task_id>/testbed/data/` and the agent may only read/write under that directory.

Your job is to decide whether the agent's trajectory correctly and completely solves the given task. DO NOT use your own factual world knowledge for verification. Observation contents may be synthetic - judge only from what the trajectory itself shows.

The solution MUST satisfy these requirements:

1. Reaches completion via finish_task:
   - The trajectory ends with a `system.finish_task` action and an observation indicating the task finished. A trajectory cut off by max-steps without `finish_task` does NOT pass.

2. Correct answer (for question-style tasks):
   - If the task asks a question (e.g., 'how many ...', 'what is the average ...'), the `answer` field of the final `finish_task` action must be the correct value, computable from values shown in earlier observations.
   - If the task does not ask for an answer, `answer` should be `"None"`.

3. Required side effects are produced:
   - When the task requires creating or modifying files (e.g., a `.docx`, `.xlsx`, `.pdf`, `.ics`, image, or email), those files must be created/written at the correct path under `tasks/<task_id>/testbed/data/` with content that satisfies the request.
   - When the task requires actions like sending an email or creating a calendar event for one or more recipients, every required recipient/event must be acted on.

4. Information correctness and grounding:
   - Values written into files, emails, or events (names, scores, dates, subjects, etc.) must come from the task description, prior observations, or the `llm`/`shell` apps. The agent must not invent values that are supposed to be extracted.
   - When extracting from a source (an excel sheet, an email body, a pdf), the agent must actually read it before writing derived data.

5. Correct app workflow:
   - Each app-specific action is executed while that app is active. A few wrong-app errors followed by `switch_app` recovery are acceptable; persistent thrashing without progress is not.

6. No unjustified shortcuts or fabrications:
   - The agent does not skip required reads (e.g., 'compute the average' without ever reading the cells).
   - The agent does not write placeholder content when the task asks for specific extracted/computed content.

7. Action format:
   - Actions are valid JSON inside `<answer> ... </answer>`. Repeated format errors that prevent progress should fail the trajectory.

8. Safety scope:
   - The agent only manipulates files under `tasks/<task_id>/testbed/data/`.

Minor inefficiencies are acceptable (extra reads, an unnecessary app switch). What matters is that the required end-state - the right files, the right answer, the `finish_task` call - is achieved on evidence visible in the trajectory.
\end{lstlisting}
\end{figure*}


% --- OfficeBench judge user prompt ---
\begin{figure*}[t]
\begin{lstlisting}[style=prettyprompt, captionpos=b, caption={OfficeBench trajectory judge -- user prompt. OfficeBench supervisors expose only \texttt{first\_name}, so the task header is shorter than AppWorld's.}, label={lst:judge_officebench_user}]
Here is the task:
<task>
My name is {first_name}.
{instruction}
</task>

Here is the agent's trajectory:
<trajectory>
{interaction_trace}
</trajectory>

Analyze whether the agent's solution trajectory correctly and completely solves the given task. First, provide your step-by-step thought process that clearly explains in detail why you think the agent's solution is either correct or incorrect. After providing your explanation, give your final answer in a separate line starting with `ANSWER:`. Your final answer must be either YES or NO. DO NOT provide anything else after the final YES/NO answer.
\end{lstlisting}
\end{figure*}
